% ======================= COMPILATION NOTICE =======================
% Before compiling, place ge2340_project.sty and all three template
% files (phase1_template.tex, phase2_template.tex, and
% contribution_template.tex) in the same directory.
% Keep the style file named ge2340_project.sty so that the templates
% can load it correctly.
% ==================================================================

\documentclass{article}

% ======================= REQUIRED STYLE =======================
\usepackage[preprint]{ge2340_project}

\usepackage[utf8]{inputenc}
\usepackage[T1]{fontenc}
\usepackage{hyperref}
\usepackage{url}
\usepackage{booktabs}
\usepackage{amsmath,amssymb}
\usepackage{graphicx}
\usepackage{microtype}
\usepackage{xcolor}
\usepackage{float}

% ======================== ASSIGNMENT CONSTRAINTS ========================
% Phase 2 is GROUP WORK. Each group has 3--4 students.
% The report must be in English and submitted as a PDF.
% The complete PDF must be 4--6 pages, including references.
% Do NOT rename, remove, or reorder the required section titles below.
% Use an appropriate number of tables and figures for qualitative and
% quantitative results. Do not paste large blocks of source code.
% Replace only the bracketed placeholders and remove instructional text before submission.
% =======================================================================

\title{Phase 2: AI Project Report\\
\large Campus Food Chatbot}

% Use the fourth author block only for a four-member group.
\author{
  \textbf{LIU, Cheuk Ying (Student ID: 58538700)}\\
  \texttt{cheuliu2-c@my.cityu.edu.hk}
  \And
  LAU, Wing Lam (Student ID: 58538853)\\
  \texttt{wllau76-c@my.cityu.edu.hk}
  \AND
  WONG, Chung Kai (Student ID: 58529125)\\
  \texttt{kaicwong68-c@my.cityu.edu.hk}
  \And
  LAN, Ka Sen (Student ID: 60166704)\\
  \texttt{kslan2-c@my.cityu.edu.hk}
}

\begin{document}
\maketitle

\noindent\textbf{Project repository (if applicable):} \url{https://github.com/[organisation-or-user]/[repository]}\\
\textbf{Project demo or supplementary link (if applicable):} \url{https://[link]}

% ======================= Required sections (do NOT rename) =======================

\section{Introduction \& Problem Recap}
Briefly recap the campus or student-life problem, target users, project scope, and intended AI task. State the implemented solution and summarize the main project outcome. Keep this section concise rather than repeating the full Phase 1 proposal.

[Replace this instruction with your report content.]

\section{Data Collection \& Methodology}
Describe how the data were obtained, created, sampled, labelled, cleaned, and split. Report relevant dataset sizes and characteristics. Explain the selected AI methodology, baselines, evaluation design, and any privacy, consent, fairness, or safety measures. Justify major methodological choices.

[Replace this instruction with your report content.]

\section{Implementation Details}
Describe the system pipeline, model/configuration, tools, training or inference procedure, and deployment or interface design. Provide enough detail for readers to understand and reproduce the workflow. Do not paste large blocks of code; use a public repository link when the project involves extensive coding.

% A compact architecture or workflow figure is encouraged.
% \begin{figure}[t]
%   \centering
%   \fbox{\rule[-0.4cm]{0cm}{3.2cm}\rule[-0.4cm]{11cm}{0cm}}
%   \caption{System architecture or end-to-end workflow (replace placeholder).}
%   \label{fig:workflow}
% \end{figure}

[Replace this instruction with your report content.]

\section{Results \& Evaluation}
Present qualitative and quantitative results using appropriate metrics, baselines, tables, figures, examples, or user feedback. Explain the evaluation protocol and analyze whether the results support the project goals. Include important negative results or failure cases rather than reporting only successful examples.

% \begin{table}[t]
%   \centering
%   \caption{Main evaluation results (replace metrics and values).}
%   \label{tab:main-results}
%   \begin{tabular}{lcc}
%     \toprule
%     Method/Setting & Metric 1 $\uparrow$ & Metric 2 $\downarrow$ \\
%     \midrule
%     [Baseline] & [\ldots] & [\ldots] \\
%     [Implemented solution] & [\ldots] & [\ldots] \\
%     \bottomrule
%   \end{tabular}
% \end{table}

[Replace this instruction with your report content.]

\section{Discussion}
\subsection{Pros and Cons of the System} 
The main strength of our system is its balance between user interaction and strict constraint enforcement. By using an LLM for intent extraction, the system can better understand the user's request flexibly and conversationally. For example, a user may enter a prompt: I want a hot meal under 50 dollars and I am allergic to peanuts. The LLM then identifies the user's budget, food preference, and allergy requirement from the input and provides a good recommendation. Besides, the system validates the recommended menu items against the PostgreSQL database before showing the results to the user. This step helps to reduce the risk of LLM hallucination and make sure the suggested items are in the database to improve user experience. However, there is a weakness that the system relies too heavily on the completeness of the database. If the canteen changes its menu and the database is not updated in a timely manner, the system may provide invalid or incorrect recommendations that do not accurately reflect the current food options. As a result, users may receive recommendations for unavailable menu items, which could negatively affect user satisfaction and trust in the system. Additionally, the system is relied on external LLM APIs which may also have potential latency of the system affected by third-party.
\subsection{Practical Benefits and Risks}
Practically, the preference data act as a reference to reduce food waste and adjust the menu to the operational company. Moreover, this system reduces the time that students and staff spend on choosing dining options effectively and provides a safety safeguard for users with a history of severe allergies. However, food safety and user trust remain major concerns based on our design. If the chatbot suggests a dish that is incorrectly tagged, it could lead to user dissatisfaction. In the worst-case scenario, this could lead to an allergic reaction, which may have serious health consequences. Therefore, it is necessary to maintain accurate and up-to-date menu data to ensure the reliability and safe deployment of the system.
\subsection{Robustness and Generalizability}
For robustness, the system is improved by validating recommendations against the structured database and applying deterministic filtering before performing the ranking process. However, the system has difficulty processing some ambiguous user requests or situations. Therefore, the satisfaction of the score is affected. Fortunately, the system is highly generalizable. The overall architecture is not restricted to a specific campus environment and logic is separated from the data source. It just needs to replace the database with another database, the same system can be deployed within closed food-service ecosystems such as in other universities, or even shopping malls.
\subsection{Team Experiences and Lessons Learned}
During the data collection and implementation stage, the team encountered several challenges that were not obvious during the initial design stage. When we were collecting data, one unexpected finding was that data preparation proved far more complex than expected. Not all menu data could be obtained through APIs. For example, some data, like AC1 canteen, could be retrieved via an API. But some data, like AC1 food cart, required taking pictures and transforming them into text. This additional data processing required more time and effort than originally planned and our schedule may be delayed. Nevertheless, this step is crucial as the quality and completeness of the collected data directly affect the accuracy and reliability of the recommendation system, so our team opened a meeting to discuss and revise the project timeline in early stage to make sure the project is smooth progress.

Another challenge during the implementation stage is setting an appropriate ranking formula for the system. Initially, the team assigned equal weight to all recommendation criteria such as price, distance, dish preferences, and food temperature. However, the test results indicated that the recommendations did not always align with user expectations. For instance, a dish might receive a high ranking because it matched food price, even if it was located very far from the user. In practice, users may prioritize convenience over a small price difference, especially when they are tired or have limited time between classes. Therefore, we tried to refine the scoring formula by assigning different weights to each criterion according to its importance. This experience highlights the difficulty of balancing diverse user preferences and demonstrates that recommendation quality depends not only on accurately capturing user intent but also on effective ranking strategy.

\section{Conclusion}
Summarize the implemented solution, the most important results, the team's unique findings, the principal limitations, and the most promising next step. Do not introduce new experiments or unsupported claims here.

[Replace this instruction with your report content.]

% ======================= References (included in page limit) =======================
% Cite sources in the text and replace the sample entry below.
\begin{thebibliography}{9}
\small
\bibitem{sample}
[Author(s)], ``[Title],'' \emph{[Venue or Source]}, [Year].
\end{thebibliography}

% Course policy reminder: LLMs may be used to refine and polish writing,
% but may not be used to generate the submission text from scratch.

\end{document}
